Language models now often read questions inside tool results, web pages, logs and emails. We test whether such a wrapper makes a model lose answers it knows, and where in the network the answer is lost. We asked four open models (Qwen3-4B, Qwen3-8B, OLMo-3-7B and Gemma-4-12B) \NItems{} short questions about facts, arithmetic and translation, kept for each model only those it answers correctly when asked plainly, and placed each inside five kinds of wrapper at three lengths. Tool results, web pages and email threads lose \LostStruct{} of known answers, depending on the model, against a noise floor of \FloorLost{}; plain prose of the same length loses about as many, and server logs lose \LostLog{}. Logs mostly change the task, so that the reply describes the log, whatever the model knows. The other wrappers mostly turn weakly held answers into wrong answers, or into claims that the text holds no answer. To find where the answer is lost, we copied internal states from the plain prompt into failed wrapped prompts. With another question's states at the question's tokens, the model gives that other question's answer in up to \SwapPeak{} of failures, as often as when the other question is asked in the wrapper. Restoring the failed question's own states helps in at most \RestorePeak{}. Restoring the last position helps only from \HandoffDepth{} of the network's depth, where the question's content reaches it. The model reads the question through the wrapper; the answer is lost later, at the position where the reply begins.
